\apendice{Documentación de usuario}

\section{Introducción}
El presente anexo tiene como finalidad proporcionar la información necesaria para que un usuario pueda instalar, ejecutar y utilizar el sistema desarrollado.

Se describen los requisitos necesarios para su funcionamiento, el proceso de instalación y las principales funcionalidades disponibles a través de la interfaz de monitorización.

\section{Requisitos de usuarios}
Para utilizar el sistema es necesario disponer de un equipo con sistema operativo Windows, Linux o macOS que tenga instalado Python 3.10 o superior.

Asimismo, es necesario disponer de las dependencias utilizadas por el proyecto, principalmente FastAPI, Streamlit y las bibliotecas auxiliares incluidas en el fichero de requisitos.

Los requisitos hardware son reducidos, ya que el sistema ha sido diseñado como un SIEM ligero. El proyecto puede ejecutarse en equipos de características modestas e incluso en plataformas de bajo consumo como Raspberry Pi.

También se requiere un navegador web moderno para acceder al dashboard de monitorización y a la documentación de la API.

\section{Instalación}
Para instalar el proyecto es necesario descargar el código fuente desde el repositorio y situarse en el directorio principal del proyecto.

Una vez descargado el proyecto, se recomienda crear un entorno virtual de Python e instalar las dependencias necesarias.

Tras la instalación de las dependencias, el sistema puede ponerse en funcionamiento ejecutando el núcleo principal del SIEM, la API REST y el dashboard de monitorización.

La aplicación web estará disponible a través del navegador y permitirá consultar toda la información generada por el sistema.

\section{Manual del usuario}
El acceso al sistema se realiza mediante el dashboard web desarrollado con Streamlit.

La página principal muestra el estado general del sistema, estadísticas de funcionamiento, anomalías detectadas y actividad reciente. Esta información permite al usuario obtener una visión global del estado de la instalación monitorizada.

La página de eventos y alertas permite consultar los registros almacenados por el sistema. El usuario puede utilizar los filtros disponibles para localizar información concreta y analizar la actividad registrada.

La página de sensores y configuración muestra el estado de los dispositivos monitorizados y permite consultar determinados parámetros de funcionamiento del sistema.

Las anomalías detectadas aparecen clasificadas según su nivel de riesgo, facilitando la identificación de aquellas situaciones que requieren una atención prioritaria.

El dashboard incorpora además diferentes gráficos y estadísticas que permiten visualizar de forma rápida la evolución de los eventos y alertas generados por el sistema.


